\clearpage
\begin{flushright}
  \fechaclase{21 de octubre de 2026}
\end{flushright}

\subsection{\texorpdfstring{\tituloejemplo{SMC--AEM para sistemas de orden superior}}{SMC-AEM para sistemas de orden superior}}

Sea el sistema no lineal SISO
\[
  \dot{X}=Ax+f(x,t)+Bu
\]
\[
  y=Cx
\]
donde
\[
  A=
  \begin{bmatrix}
    0 & 1 & 0 & \cdots & 0\\
    0 & 0 & 1 & \cdots & 0\\
    \vdots & \vdots & \vdots & \ddots & \vdots\\
    0 & 0 & 0 & \cdots & 1\\
    0 & 0 & 0 & \cdots & 0
  \end{bmatrix},
  \qquad
  B=
  \begin{bmatrix}
    0\\
    0\\
    \vdots\\
    0\\
    1
  \end{bmatrix}
\]
\[
  C=\begin{bmatrix}1 & 0 & 0 & \cdots & 0\end{bmatrix}
\]
\[
  f(x,t)=
  \begin{bmatrix}
    f_1(x,t)\\
    f_2(x,t)\\
    \vdots\\
    f_{n-1}(x,t)\\
    f_n(x,t)
  \end{bmatrix}
\]
De manera equivalente,
\begin{align*}
  \dot{x}_1 &= f_1(x,t)+x_2\\
  \dot{x}_2 &= f_2(x,t)+x_3\\
  &\ \vdots\\
  \dot{x}_{n-1} &= f_{n-1}(x,t)+x_n\\
  \dot{x}_n &= f_n(x,t)+u
\end{align*}

Las componentes $f_1,\ldots,f_{n-1}$ no están en estado de $u$, por lo que
son perturbaciones \textit{mismatched}, y $f_n$ sí lo está, por lo que es
\textit{matched}.

Se trabaja bajo las suposiciones:
\begin{enumerate}
  \item $\lVert f(x,t)\rVert_L=
  \displaystyle\sum_{i=1}^{n}\lvert f_i(x,t)\rvert\leq L$,
  con un real $L>0$.
  \item La referencia $r$ es $n$ veces diferenciable y se conoce $r^{(n)}$.
  \item El estado es medible.
\end{enumerate}

Se desea realizar seguimiento de trayectorias para la salida; los errores son:
\begin{align*}
  e_1 &= r-y=r-x_1\\
  e_2 &= \dot{r}-x_2\\
  &\ \vdots\\
  e_{n-1} &= r^{(n-2)}-x_{n-1}\\
  e_n &= r^{(n-1)}-x_n
\end{align*}
Con derivadas temporales,
\begin{align*}
  \dot{e}_1
    &=\dot{r}-\dot{x}_1
     =\dot{r}-x_2-f_1(x,t)
     =e_2-f_1(x,t)\\
  \dot{e}_2
    &=\ddot{r}-\dot{x}_2
     =\ddot{r}-x_3-f_2(x,t)
     =e_3-f_2(x,t)\\
  &\ \vdots\\
  \dot{e}_{n-1} &=e_n-f_{n-1}(x,t)\\
  \dot{e}_n
    &=r^{(n)}-f_n(x,t)-u
\end{align*}
De modo que
\[
  \dot{e}_i=e_{i+1}-f_i(x,t),
  \qquad i=1,2,\ldots,n-1
\]

Haciendo
\[
  e=\begin{bmatrix}e_1 & e_2 & \cdots & e_{n-1} & e_n\end{bmatrix}^{T}
\]
tenemos
\[
  \dot{e}=Ae-F-B\dot{e}_n
\]
\[
  \dot{e}=Ae-f-B\left(u-r^{(n)}\right)
\]
con
\[
  F=\begin{bmatrix}f_1 & f_2 & \cdots & f_{n-1} & 0\end{bmatrix}^{T}
\]

Proponemos la superficie de deslizamiento
\[
  \sigma=m^{T}e
\]
con
\[
  m=\begin{bmatrix}m_1 & m_2 & \cdots & m_{n-1} & 1\end{bmatrix}^{T}
\]
donde $m_i$ son variables de diseño que resultan en
\begin{align*}
  \sigma
    &=e_n+m_{n-1}e_{n-1}+\cdots+m_2e_2+m_1e_1\\
    &=e_n+\sum_{i=1}^{n-1}m_i e_i
\end{align*}

Derivando la superficie,
\[
  \dot{\sigma}=m^{T}\left[Ae-f-B\left(u-r^{(n)}\right)\right]
\]
Se propone la función de Lyapunov $V=\frac{1}{2}\sigma^2$ y se deriva:
\[
  \dot{V}=\sigma\dot{\sigma}
  =\sigma m^{T}\left[Ae-f-B\left(u-r^{(n)}\right)\right]
\]

Proponemos
\[
  u=r^{(n)}+m^{T}Ae+K\operatorname{sign}(\sigma)
\]
resultando
\[
  \dot{V}=\sigma\left[-m^{T}f-K\operatorname{sign}(\sigma)\right]
\]
Usando las suposiciones,
\[
  \lVert m^{T}f\rVert
  \leq\lVert m\rVert_L+\lVert f\rVert_L
  \leq m_n+L\leq M
\]
con $\lVert m\rVert\leq m_n$ para un real $m_n>0$, entonces
\begin{align*}
  \dot{V}
    &\leq\left(\lVert m\rVert_L+\lVert f\rVert_L\right)
      \lvert\sigma\rvert-K\lvert\sigma\rvert\\
    &\leq m\lvert\sigma\rvert-K\lvert\sigma\rvert\\
    &\leq-(K-m)\lvert\sigma\rvert
\end{align*}

Para $K>0$ tenemos $\dot{V}<0$ y el tiempo de alcance cumple la cota
\[
  t_r\leq\frac{\lvert\sigma(0)\rvert}{K-m}
\]
Una vez cumplido este tiempo, $\sigma=0$ y entonces
\[
  \sigma=e_n+\sum_{i=1}^{n-1}m_i e_i=0
\]
\[
  e_n=-\sum_{i=1}^{n-1}m_i e_i
\]